% Problems and Solutions (BGU \S 10.7.2.8).
% V6: readability pass over the V5 entries (same five problems, same
% symptom / fix / guard structure, simpler sentences). Leads with the three
% campaign-defining methodological problems, then hardware bring-up and the
% reproducibility exposure.
\section{Problems and Solutions}
\label{sec:problems}

This section describes the substantive engineering and scientific problems
met during the 402-cell campaign. Each entry pairs the observed symptom with
the corrective action and, where relevant, the automated guard that keeps the
problem from coming back.

\subsection{Native-resolution refactor and its rollback}
\textbf{Symptom.}\quad An intermediate design (the ``V3'' refactor) trained
each backbone at its dataset-native resolution --- $32 \times 32$ for the
CIFAR-10 inputs to TransNeXt, $28 \rightarrow 32$ for MNIST --- on the
argument that upsampling to $224 \times 224$ wastes computation and adds
interpolation artifacts. The argument sounds efficient but breaks the one
invariant the study depends on: with different tensor shapes per backbone,
any measured cross-model gap mixes the architecture with the input size it
happened to receive. It also threw away the ImageNet-pretrained
receptive-field hierarchy of the two CNNs (calibrated for $224 \times 224$)
and moved TransNeXt off the 224-pretrained distribution of its checkpoint.

\textbf{Fix.}\quad Rolled back to a single uniform $224 \times 224$ protocol:
every input is upsampled by the data pipeline before any model sees it, so
cross-backbone deltas are attributable to architecture alone. The decision is
recorded as a frozen invariant in the technical specification
(Section~\ref{sec:spec}), and the data module asserts the $224 \times 224$
output shape at construction time, so the native-resolution path cannot be
reintroduced silently.

\subsection{Hyperparameter tuning at campaign scale: frozen Optuna winners}
\textbf{Symptom.}\quad Tuning each of the 402 cells separately was both
computationally infeasible and scientifically wrong: per-level tuning would
let the optimizer compensate for each degradation level individually,
confounding the robustness signal with search noise, and a blind wide search
risked overfitting the validation split.

\textbf{Fix.}\quad The search was constrained twice. The space was narrowed
to the paper-derived priors per architecture (at most one decade around the
published value), and each (model, dataset) pair was tuned \emph{once}, at
the L3 operating point, with the winner \emph{frozen} for the entire sweep.
Phases that add regularization (D) or simplify the protocol (B2) layer on
top of the frozen winners rather than re-tuning, so every comparison happens
at one documented configuration. The winners live under
\texttt{artifacts/best\_hparams/} and round-trip through each cell's
\texttt{metrics.json}, making every run auditable.

\subsection{Blur that was invisible at \texorpdfstring{$224 \times 224$}{224x224}}
\textbf{Symptom.}\quad The single-axis blur sweep produced almost no accuracy
drop between L1 and L5 (under 1~\Mpp{} on DenseNet121) --- suspicious for the
supposedly second-worst degradation. The cause: the blur kernels
(K $= 3 \ldots 13$, $\sigma = 0.8 \ldots 2.3$) had been calibrated for native
$32 \times 32$ imagery. On the upsampled $224 \times 224$ tensor they covered
only 1--6\% of the image width --- the ``blur'' axis was not actually
blurring.

\textbf{Fix.}\quad The schedule was rescaled to pixel-domain values at
$224 \times 224$ (K $= 13 \ldots 91$, $\sigma = 2.5 \ldots 18$), tied by the
rule $K = 2\lceil 2.5\sigma \rceil + 1$ so the kernel always captures the
$2.5\sigma$ Gaussian footprint. A unit test asserts that L5 blur costs at
least 5~dB of PSNR against L1, guarding any future schedule against the same
scale mistake.

\subsection{Vision-transformer bring-up on Blackwell hardware}
\textbf{Symptom.}\quad TransNeXt training crashed on the RTX 5070 (Blackwell,
sm\_120) with an opaque CUDA kernel-launch error in the first attention
block: the PyTorch stack at project start shipped kernels only up to sm\_90,
and the TransNeXt \texttt{swattention} extension had no sm\_120 binary.

\textbf{Fix.}\quad TransNeXt was quarantined until the upstream package was
rebuilt with sm\_120 support in the PyTorch Blackwell release; the quarantine
predicate is now a permanent no-op but stays in code for future GPU
substitutions. A startup self-test runs a one-iteration forward pass on the
configured GPU before the campaign driver dispatches its first cell, so a
compute-capability mismatch surfaces as a structured error instead of a raw
CUDA stack trace mid-campaign.

\subsection{Single-seed reproducibility exposure}
\textbf{Symptom.}\quad The L3 headline cells were originally single-seed
point estimates at seed 42. Without a variance bound, a reviewer could not
tell whether a small inter-model gap was a real effect or a lucky seed.

\textbf{Fix.}\quad A 48-replicate audit re-trained all 24 L3 headline cells
at seeds 43 and 44. Every audited cell shows $\sigma < 1.2$~\Mpp{} (worst:
the TransNeXt-tiny CIFAR-10 Phase~D T3 cell at 1.18), and no cross-model
conclusion reverses under any seed. Schema version 3 of
\texttt{Final\_Exp.json} carries the mean-$\pm\sigma$ pair on every audited
cell, so the reproducibility floor is visible without re-running the audit.
